\documentclass[a4paper,11pt]{article}

\usepackage{amsmath,amssymb,mathtools}
\usepackage{booktabs}
\usepackage{listings}
\usepackage{xcolor}
\usepackage{hyperref}
\usepackage{geometry}
\geometry{margin=25mm}

\definecolor{codegray}{gray}{0.95}
\lstset{
  basicstyle=\ttfamily\small,
  backgroundcolor=\color{codegray},
  frame=single,
  columns=fullflexible,
  breaklines=true,
  keepspaces=true
}

\title{Computational Experiments with Stein Factorization in Macaulay2}
\author{Takehiko Yasuda}
\date{July 2026}

\begin{document}

\maketitle

\begin{abstract}
We implement a prototype of a Stein factorization algorithm in Macaulay2 and test it on
several explicit projective maps. The implementation follows the bigraded Hom-module
construction in the author's algorithm:
\[
 C=\operatorname{Hom}_R(R_{\geq r},R)_{(0,\geq 0)}.
\]
The module is used to reconstruct a graded coordinate algebra for the Stein intermediate
scheme and, where possible, the graph of the connected-fiber morphism. This note is not
intended as a formal verification of the algorithm; rather, it records a reproducible
computational experiment that demonstrates the practical feasibility of the construction
on concrete examples. All calculations are included in the repository and can be reproduced
with Macaulay2 1.24.11 or later. The code was developed with the assistance of an AI
system; the mathematical setup, examples, expected structures, and verification criteria
were selected and reviewed by the author.
\end{abstract}

\section{Purpose and background}

For a projective morphism
\[
 f:Y\longrightarrow X,
\]
a Stein factorization is a decomposition
\[
 Y\xrightarrow{h} Z\xrightarrow{g} X
\]
in which $h$ has connected fibers and $g$ is finite. The purpose of this project is to
turn an algorithm described by Takehiko Yasuda\cite{yasuda} into an executable Macaulay2
prototype and to see how it behaves on concrete examples from projective geometry.

The input is the graph of $f$ represented as a bigraded projective scheme. If
\[
 R=S/I_{\Gamma_f},
\]
then the variables in the source block have degree $(1,0)$ and those in the target block
have degree $(0,1)$. The two gradings make it possible to separate the source direction
from the target direction during the Hom-module computation.

\section{The computational construction}

For a sufficiently large bidegree $r=(r_1,r_2)$, the implementation computes
\[
 C=\operatorname{Hom}_R(R_{\geq r},R)_{(0,\geq0)}.
\]
The truncation bound is obtained from shifts in a bounded bigraded free resolution of
$R$ over the ambient polynomial ring. This is the computational counterpart of the
bound in Corollary 4.3 of \cite{yasuda}.

The main public functions are:

\begin{lstlisting}
data  = steinHomData(S, graphIdeal, d1, d2, c1, c2)
cData = steinCoordinateAlgebra(data, gammaIndex, baseImages)
gData = directSteinGraph(data, cData)
\end{lstlisting}

The first function computes the truncated Hom module and records the bound. The second
evaluates selected Hom generators at a nonzero element $\gamma$ and computes the defining
relations of the resulting coordinate algebra. The third computes a graph closure from
the evaluated generators in a localization. For examples where the minimal resolution
is too expensive, the experimental entry point
\texttt{steinHomDataAtBound} accepts a supplied bound. Such a computation is explicitly
marked as non-certified in the output.

\section{Environment and reproducibility}

The computations were run with Macaulay2 1.24.11 on macOS (Apple Silicon, 2023). The relevant
bundled packages include \texttt{Truncations}, \texttt{MinimalPrimes}, and \texttt{Saturation}.
Standard examples complete in under 5 seconds each; the full suite including intermediate
examples runs in approximately 30--60 seconds total.

From the project root, the standard suite is run by:

\begin{verbatim}
./run-tests.sh
\end{verbatim}

The slower twisted-cubic benchmark is run separately:

\begin{verbatim}
M2 --no-readline --stop -q tests/blowup-twisted-cubic.m2
\end{verbatim}

The complete input files are included in the repository. The excerpts below are chosen
to show the mathematical input and the checks, without reproducing every elimination
polynomial or the full Macaulay2 startup log.

\section{Representative Macaulay2 calculations}

\subsection{A quadratic map of $\mathbb{P}^1$}

The first test is the finite map $[s:t]\mapsto[s^2:t^2]$. Its graph is described by
$y_0^2x_1-y_1^2x_0$. The following excerpt shows the input and the main output.

\begin{verbatim}
i = ideal(y0^2*x1-y1^2*x0);
d = steinHomData(S,i,1,1,2,2);
c = steinCoordinateAlgebra(d,0,{x0,x1});

d#"bound"
o = {1, 0}

d#"steinGeneratorDegrees"
o = {{0, 0}, {0, 1}}

c#"definingIdeal"
o = ideal (X0*X1-z^2)

isPrime (directSteinGraph(d,c)#"graphIdeal")
o = true
\end{verbatim}

Thus the Stein intermediate is the source $\mathbb{P}^1$, displayed as the conic
$X_0X_1-z^2=0$ over the target coordinate ring.

\subsection{A Mori fiber space with a finite part}

The next example is the composition
\[
 \mathbb{P}^1\times\mathbb{P}^2\longrightarrow\mathbb{P}^2
 \xrightarrow{[x_0:x_1:x_2]\mapsto[x_0^2:x_1^2:x_2^2]}
 \mathbb{P}^2.
\]
The source projection has connected $\mathbb{P}^1$-fibers, while the second map is
finite. The relevant part of the test input and output is:

\begin{verbatim}
igraph = kernel graphParametrization;
d = steinHomData(sgraph,igraph,5,2,6,3);
c = steinCoordinateAlgebra(d,0,{y0,y1,y2});

numgens igraph, codim igraph
o = (12, 4)

d#"bound"
o = {2, 0}

dim(c#"ring")
o = 3

hilbertFunction(1,c#"ring"), hilbertFunction(2,c#"ring")
o = (6, 15)

isPrime(c#"definingIdeal")
o = true

fiberComponents = minimalPrimes ifiber;
#fiberComponents
o = 4
\end{verbatim}

Each of the four fiber components is a copy of $\mathbb{P}^1$. The computed Stein
intermediate is the degree-two Veronese model of $\mathbb{P}^2$, and the connected part
is the projection to $\mathbb{P}^2$.

\subsection{The blow-up of a line in $\mathbb{P}^3$}

We next take the blow-down $\operatorname{Bl}_L(\mathbb{P}^3)\to\mathbb{P}^3$ along a
line $L$, followed by the coordinatewise square map on $\mathbb{P}^3$. The main numerical
checks are:

\begin{verbatim}
numgens igraph, codim igraph
o = (26, 7)

d = steinHomData(sgraph,igraph,7,3,8,4);
d#"bound", #d#"steinGeneratorIndices"
o = ({2, 0}, 8)

dim(c#"ring")
o = 4

hilbertFunction(1,c#"ring"), hilbertFunction(2,c#"ring")
o = (10, 35)

isPrime c#"definingIdeal"
o = true

degree(ssource/ie)
o = 2

#fiberPrimes, all(fiberPrimes,p -> degree(rfiber/p) == 1)
o = (8, true)
\end{verbatim}

The Stein ring has the Hilbert values of the second Veronese ring of $\mathbb{P}^3$.
The exceptional divisor has degree two and is the expected Segre quadric
$\mathbb{P}^1\times\mathbb{P}^1$. A representative fiber consists of eight reduced point
components.

\subsection{The twisted-cubic blow-up}

The final example is the blow-up of $\mathbb{P}^3$ along the twisted cubic, followed by
the coordinatewise square map. It is deliberately retained as a heavier benchmark.
The full minimal resolution needed for the automatic Corollary 4.3 bound is expensive
(typically 10--30 seconds of computation on a 2023 MacBook Pro), so the test supplies
$r=(2,0)$ and records that this bound is experimental.

\begin{verbatim}
reesI = kernel map(targetRees,trees,{...});
numgens reesI, codim reesI
o = (3, 2)

igraph = saturate(igraphRaw,ideal matrix{flatten rows});
numgens igraph, codim igraph, isPrime igraph
o = (70, 11, true)

d = steinHomDataAtBound(sgraph,igraph,{2,0});
d#"steinGeneratorDegrees"
o = {{0,0},{0,1},{0,1},{0,1},{0,1},{0,1},{0,1},{0,2}}

hilbertFunction(1,c#"ring"), hilbertFunction(2,c#"ring")
o = (10, 35)

d#"certifiedBound"
o = false
\end{verbatim}

The resulting ring again has the Hilbert values of the second Veronese of $\mathbb{P}^3$.
One degree-two module generator is algebraically redundant, so automatic graph
construction is not yet completed for this example. This is a useful illustration of the
difference between a successful numerical stabilization and a certified bound.

\section{Summary of results}

\begin{table}[ht]
\centering
\caption{Main computational results}
\begin{tabular}{lllll}
\toprule
Example & bound & $\dim C$ & selected Hilbert data & graph\\
\midrule
$\mathbb{P}^1$ quadratic map & $(1,0)$ & 2 & $H(1)=3$ & prime\\
$\mathbb{P}^1$ cubic map & $(2,0)$ & 2 & twisted cubic & prime\\
$\mathbb{P}^1\times\mathbb{P}^2$ & $(2,0)$ & 3 & $H(1)=6,H(2)=15$ & prime\\
$\operatorname{Bl}_L(\mathbb{P}^3)$ & $(2,0)$ & 4 & $H(1)=10,H(2)=35$ & prime\\
Twisted-cubic blow-up & $(2,0)^*$ & 4 & $H(1)=10,H(2)=35$ & partial\\
\bottomrule
\end{tabular}
\end{table}

The star indicates a supplied experimental bound rather than a bound computed from the
full minimal resolution. Across these examples, the implementation reproduces the
expected intermediate objects for finite maps, positive-dimensional connected fibers,
and a divisorial contraction. The graph ideals produced by the direct localization
strategy are prime in the completed examples.

\section{Limitations}

This is a research prototype rather than a general-purpose Macaulay2 package. The main
limitations are:

\begin{itemize}
  \item deep minimal free resolutions can dominate the exact-bound computation;
  \item finite module generators may be redundant as algebra generators;
  \item heuristic stabilization gives evidence but does not prove the explicit bound;
  \item the component-selection part of the full algorithm is not automated for every
        possible input.
\end{itemize}

The calculations should therefore be read as reproducible computational evidence, not as
a replacement for the mathematical hypotheses or proofs in the source paper.

\section{Use of AI in the implementation}

An AI system was used to help draft and revise Macaulay2 code, test cases, and explanatory
documentation. The mathematical formulation, choice of examples, expected geometric
interpretations, and verification criteria were supplied and reviewed by the author.
The generated code was executed in Macaulay2, and the outputs were checked using defining
ideals, Hilbert functions, primality tests, fiber decompositions, and composition checks.
In particular, computations marked as non-certified were not presented as proofs.

\section{Conclusion}

This project shows that the bigraded Hom-module construction can be used in a practical
Macaulay2 workflow on several nontrivial examples. The most useful outcome is not a large
collection of examples, but a transparent record linking the mathematical construction,
the source code, the actual computer input, and the numerical output. The repository is
therefore intended as a public technical note and reproducible experiment rather than as
a claim of a fully general implementation.

\begin{thebibliography}{9}

\bibitem{hartshorne}
Robin Hartshorne,
\emph{Algebraic Geometry},
Graduate Texts in Mathematics, Vol. 52, Springer-Verlag, 1977.
(Chapter II, \S3 on separated and proper morphisms; Chapter III on coherent sheaves.)

\bibitem{beauville}
Arnaud Beauville,
\emph{Variétés Kähleriennes dont la première classe de Chern est nulle},
Journal of Differential Geometry, Vol. 18, No. 4, pp. 755--782, 1983.

\bibitem{yasuda}
Takehiko Yasuda,
\emph{An algorithm for the minimal model program in dimension three},
arXiv:2603.13703v2.
\url{https://arxiv.org/abs/2603.13703}

\bibitem{m2}
Daniel R. Grayson and Michael E. Stillman,
\emph{Macaulay2, a software system for research in algebraic geometry}.
\url{https://math.uiuc.edu/Macaulay2/}

\end{thebibliography}

\end{document}
